
\chapter{מבוא לתיאום כלים מרובים בסוכני \textenglish{LLM}}

\section{האבולוציה של סוכני \textenglish{AI}}

בתחילת שנות ה-\textenglish{2020}, מודלי שפה גדולים (\textenglish{LLM}: \textenglish{Large Language Models}) \cite{brown2020language} היו בעיקרם כלים ליצירת טקסט ותשובה לשאלות. מהפכת \textenglish{Transformer} \cite{vaswani2017attention} פתחה אפשרות לסוכנים הפועלים בפידבק לולאה (\textenglish{agentic loop}): הסוכן מקבל משימה, בוחר כלי, מקבל תשובה ומחליט על הצעד הבא.

\textenglish{ReAct} \cite{yao2022react} הייתה מסגרת הפיילוט ב-\textenglish{2022}: השילוב של \textenglish{Reasoning} (\textenglish{CoT}: \textenglish{Chain-of-Thought}) עם \textenglish{Acting} (קריאה לכלים חיצוניים) אפשר פתרון מרשים של בעיות מורכבות הדורשות מידע עדכני. מאז, האבולוציה הייתה מהירה: מסוכן בודד עם כלי אחד לסביבות מרובות-סוכנים עם מאות כלים.

\section{הגדרת תיאום כלים}

תיאום כלים (\textenglish{tool orchestration}) מוגדר כתהליך של בחירה, שרשור וניהול קריאות לכלים חיצוניים על ידי סוכן-\textenglish{LLM}. מימד המורכבות הגובר כולל:
\begin{itemize}
\item בחירה דינמית מתוך מאות כלים (\textenglish{tool selection})
\item ביצוע מקבילי של כלים בלתי-תלויים (\textenglish{parallel tool calls})
\item ניהול שגיאות ו\textenglish{retry} אוטומטי
\item עלות אסימון (\textenglish{token economics}) \cite{segal2024orchestration}
\end{itemize}
